\section{丛对及其马斯洛夫指标}
\label{sec:bundle-pair-maslov-index}

设 $\Sigma$ 是带边界 $\partial\Sigma$ 的定向紧曲面。以 $g$ 表示
$\Sigma$ 的亏格，以 $h$ 表示 $\partial\Sigma$ 的连通分支数。

\subsection{复向量丛情形}
\label{subsec:complex-vector-bundle-pairs}

考虑复向量丛 $E\to\Sigma$。注意，若 $\partial\Sigma\ne\varnothing$，则
$\Sigma$ 上的任意复向量丛 $E\to\Sigma$ 在拓扑上都是平凡的。

\stfhSubIndex{I048}{Congdui}{丛对（bundle pair）}{Fu}{复丛对（complex bundle pair）}
\begin{definition}
\label{def:complex-bundle-pair}
一个\termComplexBundlePair{} $(E,\lambda)$ 是复向量丛 $E\to\Sigma$ 连同
$E|_{\partial\Sigma}$ 的一个实子丛 $\lambda\to\partial\Sigma$，并配备同构
\begin{equation}
  E|_{\partial\Sigma}\cong\lambda\otimes\mathbb C.
\label{eq:complex-bundle-boundary-complexification}
\end{equation}
等价地，复丛对是复向量丛 $E\to\Sigma$ 连同
$E_{\mathbb R}|_{\partial\Sigma}\to\partial\Sigma$ 的一个\termTotallyRealSubbundle{}
$\lambda\subset E_{\mathbb R}|_{\partial\Sigma}$；
这里 $E_{\mathbb R}$ 是丛 $E$ 的实化。
\end{definition}

固定一个\termTrivialization{}
\[
  \Phi:E\longrightarrow\Sigma\times\mathbb C^n,
\]
并设 $R_1,\ldots,R_h$ 是 $\partial\Sigma$ 的连通分支，各自带有从
$\Sigma$ 诱导的\termBoundaryOrientation{}。由给定的同构
$E|_{\partial\Sigma}\cong\lambda\otimes\mathbb C$，
$\Phi(\lambda|_{R_i})$ 给出环路
\[
  \gamma_{\Phi,i}:S^1\longrightarrow\mathcal R(n).
\]
令 $\mu(\Phi,R_i)=\mu(\gamma_{\Phi,i})$，便得到下面的命题。

\begin{proposition}
\label{prop:complex-bundle-pair-index-trivialization}
设 $(E,\lambda)$ 是定向紧曲面 $\Sigma$ 上秩为 $n$ 的复丛对。则
\[
  \sum_{i=1}^h\mu(\Phi,R_i)
\]
不依赖于平凡化 $\Phi:E\to\Sigma\times\mathbb C^n$ 的选取。\footnote{译者注：原书把求和指标写成 $j$，被求和项却仍为 $R_i$；译文统一为 $i$。}
\end{proposition}

\begin{proof}
设 $\Phi_1,\Phi_2$ 是 $E$ 的两个平凡化。于是
\[
  \Phi_2\circ\Phi_1^{-1}:
  \Sigma\times\mathbb C^n\longrightarrow\Sigma\times\mathbb C^n
\]
由 $(x,v)\mapsto(x,g(x)v)$ 给出，其中
\[
  g:\Sigma\longrightarrow GL(n,\mathbb C).
\]
由 $\gamma_{\Phi,i}$ 的定义\footnote{译者注：原书把第一条矩阵恒等式的左端写成
$g\overline g^{-1}\cdot B\circ\gamma_{\Phi_1,i}$，并在随后的度数公式中省去
$B\circ\gamma_{\Phi_j,i}$ 前的 $\det$，还把每个分支上的限制 $|_{R_i}$
写成 $|_{\partial\Sigma}$。译文依 $B(A\mathbb R^n)=A\overline A^{-1}$
及本节对 $\mu(\Phi,R_i)$ 的定义补正这些对象。}，
\[
  g\,(B\circ\gamma_{\Phi_1,i})\,\overline g^{-1}
  =B\circ\gamma_{\Phi_2,i}.
\]
因此
\[
\begin{aligned}
  \deg(\det\circ B\circ\gamma_{\Phi_2,i})
  ={}&
  \deg\left(\left.\det(g\overline g^{-1})\right|_{R_i}\right)
  +\deg(\det\circ B\circ\gamma_{\Phi_1,i}),
\end{aligned}
\]
亦即
\[
  \mu(\Phi_2,R_i)-\mu(\Phi_1,R_i)
  =\deg\left(\left.\det(g\overline g^{-1})\right|_{R_i}\right).
\]
另一方面，因为 $\partial\Sigma=\coprod_iR_i$，而映射
\[
  \left.\det(g\overline g^{-1})\right|_{\partial\Sigma}:
  \partial\Sigma\longrightarrow S^1
\]
显然可以延拓到 $\Sigma$，所以由度数的\termCobordismInvariance{}，
\[
  \sum_{i=1}^h
  \deg\left(\left.\det(g\overline g^{-1})\right|_{R_i}\right)=0.
\]
证明完成。
\end{proof}

\cref{prop:complex-bundle-pair-index-trivialization}说明下面的定义是良定义的。

\stfhIndex{I124}{Masiluofuzhibiao}{马斯洛夫指标（Maslov index）}
\begin{definition}
\label{def:complex-bundle-pair-maslov-index}
\termBundlePair{} $(E,\lambda)$ 的\termMaslovIndex{}定义为
\begin{equation}
  \mu(E,\lambda)=\sum_{i=1}^h\mu(\Phi,R_i),
\label{eq:complex-bundle-pair-maslov-index}
\end{equation}
其中 $\Phi:E\to\Sigma\times\mathbb C^n$ 是某个（因而任意）平凡化。
\end{definition}

\subsection{辛向量丛情形}
\label{subsec:symplectic-vector-bundle-pairs}

考虑一个\termSymplecticVectorBundle{} $E\to\Sigma$，也就是说，假设每条纤维
$E_x$ 都带有一个随 $x\in\Sigma$ 光滑变化的辛内积 $\omega_x$。
同样，若 $h>0$，则 $\Sigma$ 上的任意辛向量丛都是辛平凡的。固定一个\termSymplecticTrivialization{}
\[
  \Psi:E\longrightarrow\Sigma\times(\mathbb R^{2n},\omega_0).
\]

\stfhSubIndex{I049}{Congdui}{丛对（bundle pair）}{Xin}{辛丛对（symplectic bundle pair）}
\begin{definition}
\label{def:symplectic-bundle-pair}
一个\termSymplecticBundlePair{}是对子 $(E,\lambda)$，其中
$E\to\Sigma$ 是辛向量丛，而 $\lambda\to\partial\Sigma$ 是
$E|_{\partial\Sigma}$ 的一个\termLagrangianSubbundle{}。
\end{definition}

回忆，$\mathbb C^n$ 的任意拉格朗日子空间之\termComplexification{}都典范同构于
$\mathbb C^n$。把这与\cref{def:complex-bundle-pair}比较。与复情形一样，限制
$\Psi(\lambda|_{R_i})$ 给出环路
\[
  \gamma_{\Psi,i}:S^1\longrightarrow\Lambda(n).
\]
记 $\mu(\Psi,R_i)=\mu(\gamma_{\Psi,i})$。于是有下述命题。

\begin{proposition}
\label{prop:symplectic-bundle-pair-index-trivialization}
设 $(E,\lambda)$ 是定向紧曲面 $\Sigma$ 上的辛丛对。则
\[
  \sum_{i=1}^h\mu(\Psi,R_i)
\]
不依赖于辛平凡化
$\Psi:E\to\Sigma\times(\mathbb R^{2n},\omega_0)$ 的选取。\footnote{译者注：原书再次把求和指标写成 $j$、被求和项写成 $R_i$；并把平凡化的靶写成 $E\times\mathbb C^n$。译文分别统一为 $i$ 及本小节已经固定的 $\Sigma\times(\mathbb R^{2n},\omega_0)$。}
\end{proposition}

\begin{proof}
设
\[
  \Psi_1,\Psi_2:
  E\longrightarrow\Sigma\times(\mathbb R^{2n},\omega_0)
\]
是两个辛平凡化\footnote{译者注：原书在本证明中把两个辛平凡化的靶及过渡映射的定义域、值域都只写成
$(\mathbb R^{2n},\omega_0)$，但下一行的赋值仍使用基点 $x$。译文补回
$\Sigma\times(\mathbb R^{2n},\omega_0)$。}。映射
\[
  \Psi_2\circ\Psi_1^{-1}:
  \Sigma\times(\mathbb R^{2n},\omega_0)
  \longrightarrow
  \Sigma\times(\mathbb R^{2n},\omega_0)
\]
由 $(x,v)\mapsto(x,g(x)v)$ 给出，其中
\[
  g:\Sigma\longrightarrow\operatorname{Sp}(2n,\mathbb R).
\]
把 $\operatorname{Sp}(2n,\mathbb R)$ 在 $\Lambda(n)$ 上的自然作用写成
$g\cdot V$。于是
\[
  g(x)\gamma_{\Psi_1,i}(x)=\gamma_{\Psi_2,i}(x),
  \qquad x\in R_i.
\]

因为 $\mathbb C^n$ 的任意拉格朗日子空间关于标准复乘法都是全实的，而且\termGeneralizedMaslovIndex{}在 $\mathcal R(n)$ 中的同伦下不变，所以这个恒等式给出
\[
  \mu(\Psi_2,R_i)=\mu(\Psi_1,R_i)+2\deg(g|_{R_i}),
\]
其中 $\deg(g|_{R_i})$ 是第 2.2.2 节中定义的环路
\[
  g|_{R_i}:S^1\longrightarrow\operatorname{Sp}(2n,\mathbb R)
\]
的度数。因为
\[
  g|_{\partial\Sigma}=\coprod_{i=1}^h g|_{R_i},
\]
再次得到
\[
  \sum_{i=1}^h\deg(g|_{R_i})=0,
\]
证明即告完成。
\end{proof}

这个命题使我们能够定义辛丛对 $(E,\lambda)$ 的马斯洛夫指标
$\mu(E,\lambda)$。

\begin{definition}
\label{def:symplectic-bundle-pair-maslov-index}
辛丛对 $(E,\lambda)$ 的马斯洛夫指标定义为
\begin{equation}
  \mu(E,\lambda)=\sum_{i=1}^h\mu(\Psi,R_i),
\label{eq:symplectic-bundle-pair-maslov-index}
\end{equation}
其中 $\Psi:E\to\Sigma\times(\mathbb R^{2n},\omega_0)$ 是某个（因而任意）
辛平凡化。\footnote{译者注：原书此处沿用复丛对定义中的 $\Phi$ 与
$\Sigma\times\mathbb C^n$，但本小节此前始终以 $\Psi$ 表示取值于
$\Sigma\times(\mathbb R^{2n},\omega_0)$ 的辛平凡化；译文依此统一。}
\end{definition}

不难证明 $\mu(E,\lambda)$ 的\termHomotopyInvariance{}。事实上，它满足一个更强的不变性；现在来表述这一点。

设 $(E_1,\lambda_1)$、$(E_2,\lambda_2)$ 分别是 $\Sigma_1$、$\Sigma_2$
上的辛丛对。\footnote{译者注：原书此处先写 $(E_0,\lambda_0)$、
$(E_1,\lambda_1)$，随后在命题中混用下标 $0,1,2$ 以及未定义的
$\Sigma_0,E_2,\lambda_2$。译文依命题结论与配边的两个边界分支统一为下标
$1,2$。}

\begin{proposition}
\label{prop:maslov-cobordism-invariance}
假设实向量丛
\[
  \lambda_1|_{\partial\Sigma_1}
  \sqcup
  \lambda_2|_{\partial\Sigma_2}
\]
可以延拓为\termOrientedCobordism{} $\Sigma_{12}$ 上的实向量丛
$\lambda_{12}\to\Sigma_{12}$，其中
\[
  \partial\Sigma_{12}
  =-\partial\Sigma_1\sqcup\partial\Sigma_2,
\]
并且 $\lambda_{12}$ 嵌入辛向量丛 $E_{12}\to\Sigma_{12}$，成为一个拉格朗日子丛；其在边界 $\partial\Sigma_{12}$ 上限制为丛对
\[
  (E_1,\lambda_1)|_{\partial\Sigma_1}
  \sqcup
  (E_2,\lambda_2)|_{\partial\Sigma_2}.
\]
则
\[
  \mu(E_1,\lambda_1)=\mu(E_2,\lambda_2).
\]
\end{proposition}

\begin{proof}
任选 $E_{12}\to\Sigma_{12}$ 的一个平凡化
\[
  \Phi_{12}:E_{12}
  \longrightarrow
  \Sigma_{12}\times(\mathbb R^{2n},\omega_0).
\]
它显然定义了 $E_1$ 和 $E_2$ 的平凡化。在这个平凡化中计算
$\mu(E_i,\lambda_i)$（$i=1,2$）以及 $\mu(E_{12},\lambda_{12})$，结论随即成立。
\end{proof}

我们把这个命题的陈述称为 $\mu(E,\lambda)$ 的\termCobordismInvariance{}。
显然，这个命题蕴含 $\mu(E,\lambda)$ 的同伦不变性。

现在设 $L\subset(M,\omega)$ 是拉格朗日子流形，并考虑光滑映射
\[
  w:(\Sigma,\partial\Sigma)\longrightarrow(M,L).
\]
以 $\partial w:\partial\Sigma\to L$ 表示限制 $w|_{\partial\Sigma}$。于是
$(w,\partial w)$ 诱导一个典范丛对
\[
  E=w^*TM,
  \qquad
  \lambda=(\partial w)^*TL.
\]

\begin{definition}
\label{def:maslov-index-map}
设 $w$ 如上。$w$ 的马斯洛夫指标记为 $\mu_L(w)$，定义为
\[
  \mu_L(w)=\mu\bigl(w^*TM,(\partial w)^*TL\bigr).
\]
\end{definition}

$\mu(E,\lambda)$ 的同伦不变性立即推出映射 $w$ 的 $\mu_L(w)$ 具有同伦不变性。
